\renewcommand\footerfilename{Typography.tex}

\section{Typography} \label{sec:Typography}

The typographic conventions used in these notes. PyRAF uses single letter
commands and permits a minimal-abbreviation of commands. To be clear about
what needs to be typed, these keys are typeset clearly with \LaTeX.

The document uses \dhl{colors} to make \dhl{interesting} things stand
out from the text.

For example: \llbox{k} means hit the ``k'' on the keyboard. 

\mousebox{mouse-drag left-mouse-1} means to use a mouse drag
event. There are no mouse-drag events in PyRAF, but they do occur in
ds9, Demetra and other ancillary programs of mild disiterest.


\subsection*{Provided Examples:}

Examples, in general, state:
\vspace{-.15cm}
\begin{enumerate}\addtolength{\itemsep}{-0.5\baselineskip}
%\setcounter{enumi}{N}
   \item   \dhl{Goal} what we want to you to practice
   \item   \dhl{example} the gist of the example
   \item   \dhl{Prose}  describes a 'thought model' or a statement about takeaways from the example
   \item   \dhl{Exercise} - The exercise's steps
   \item   \dhl{Take Away} - emphasizes what you want to stick to your brain.
\end{enumerate}

Not all parts may appear. 

\begin{quote}
\goal Demonstrate a goal. Yes this example is a goal.\\ \example Give
an example of some syntax.\\ \prose Hey, put a linguistic pattern into
your ``thought process'' for that syntax.\\ \exercise[-a] Pick up your
pencil and start your notebook! \\ \exercise[-b] Time to stand for a
minute. \\ \takeaway Tie the thought process to a concrete goal and
example and begin backing up, documenting and checking your work.
\end{quote}

\subsection{Margin Annotations}

There are places in the notes where margin annotations (glosses) are
used to draw attention to something important. In some cases this will
be a note to the author to add/clarify a section of text. Usually more
information may be found in the End Notes section of the document.
\ltodo{Example of Margin Note}{Added a note to demonstrate the note
  process.}

\subsection{Index}

This document is written in \LaTeX offering coordinated footnotes,
bibliographic references, an index and a glossary.
